\chapter{Fazit}

Trotz der zahllosen Erweiterungsvorschläge im vorherigen Kapitel möchte ich hier noch einmal ein Fazit ziehen und festhalten, dass die Webapp bereits viele gut funktionierende Funktionen bereitstellt. Aufgrund der beschränkten Zeit und der notwendigen Einarbeitung in die Thematik benötigt die Webapp natürlich vor ihrem professionellem Einsatz noch einige Optimierung und Ergänzungen. Dennoch ist die Webapp in ihrem aktuellen Stand bereits in der Lage (bzw. wäre das nach Erstellen und Integrieren eines Backends), Trackingdaten zu visualisieren und zu bearbeiten sowie grundlegende, komplett frei konfigurierbare Analysen zu erstellen, womit das angestrebte Ziel erreicht wurde. Alle implementierten Funktionen funktionieren gut und sind sehr reaktiv, also ohne störende Wartezeiten benutzbar. An vielen Stellen wurde durch Animationen die Nutzerfreundlichkeit und Übersichtlichkeit weiter verbessert und eine (zumindest in meinen Augen) sehr ansprechende Optik erreicht.
\\
Deshalb sollte die Webapp idealerweise auch nicht nur in diesem Dokument begutachtet werden, sondern am besten selbst ausprobiert. Der Code und die Dokumentation zu dessen Ausführung befinden sich in einem GitLab Repository unter \url{https://gitlab.lrz.de/ga38vud/saasfrontend/}.
\\
\\
Mir hat die Entwicklung der Webapp viel Freude bereitet und ich hoffe, dass dies auch bei hoffentlich bald vorhandenen Anwendern der Fall sein wird. Eine leicht zu benutzende und frei zur Verfügung stehende Anwendung kann schließlich auch dazu betragen, Spielanalyse für mehr Teams aller möglichen Sportarten kostengünstig zu ermöglichen. So wird nicht nur der Wettbewerb fairer gestaltet, sondern auch interessierte Spieler oder Trainer können sich mit den Möglichkeiten, die Ihnen diese Technologie bietet, auseinandersetzen.